\section{Discussion: what this means for a regulator}\label{sec:discussion}

The evidence is simulated for the questions on rules and real for some of the questions on data, and the two do not overlap on accounts. We therefore separate what a regulator could act on now, what it should do to learn more before acting, and what should not be concluded.

\subsection{Findings that can guide rules now}

\paragraph{Price off-book blocks from a trailing mean.} Among the price and order rules we tested this is the only one that cuts the ring's gain through a mechanism that can be written down and checked: a reference price that averages the last $L$ prices credits a push only for increments that are recent and discounts them linearly with age (Proposition~\ref{prop:block}). In the simulated markets a one-day reference removes 87\% (index level) to 98\% (stock level) of the value of the block and leaves the ring between a tenth and a half of its total gain (Section~\ref{sec:blockcheck}); this holds for a block of the assumed size, and a ring whose gain is mostly trading profit (blocks ten times smaller) keeps more than 80\% of it; a ring can recover the value only by holding the price for about $L-a/2$ steps, which is itself the conduct to deter. Its cost to honest trading is the gap between the last price and the trailing mean on block trades, 0.2\% to 0.5\% of the price in our first preset, which we have not measured on real negotiated trades. One of the ten documented cases (PSH) was carried out on a negotiated system, and the rule applies where the ring's gain is realised off the order book; it does nothing against manipulation that profits on the book. We would pilot it first on a small set of stocks and compare negotiated-trade prices with the trailing mean to measure the cost before it binds.

\paragraph{Inspection capacity and the fine are exchangeable.} With a fine proportional to the gain, deterrence requires only that the probability of being caught exceed $1/m$ (Proposition~\ref{prop:deter}). Under the fitted power law, doubling the fine multiple from five to ten cuts the capacity needed to deter a small ring by a factor of about 2.4 (3.0 to 1.2 inspected accounts per window at 600 active accounts; Table~\ref{tab:capacity}). The Vietnamese caps (five times the gain for individuals, ten for organisations, as reported by the press) therefore matter for how many accounts the inspectorate must examine, and an increase in the cap for individuals would be worth about as much as a 2.4-fold increase in capacity under our model. The empirical counterpart is the finding of \citet{comerton2014} that the regulator's budget has a strong effect on both manipulation and detection. The arithmetic of a workload is simple: $B$ inspected accounts per window, $S$ stocks and $W$ windows per year give $BSW$ account reviews a year; with $B=3$, an illustrative 400 stocks and about 200 windows (a window is about 300 minutes), that is 240,000 automated reviews a year, of which only the highest-ranked need human attention. Because the saving from screening shrinks as the number of accounts per stock grows (Table~\ref{tab:capacity}: saving of 3.7 at 600 accounts and 2.0 at 2000), capacity should be allocated first to stocks with few active accounts, where screening is most informative.

\paragraph{Do not use order-level frictions or the band against rings.} Minimum resting times and cancel fees did not reduce spoofing, pump-and-dump or wash trading in the calibrated market and cost honest participants 7--8\% in spread and 21--24\% in volatility at the strictest settings; a larger tick widened the spread by 35\% (Figure~\ref{fig:tradeoff}). Widening the daily band to 10\% or 15\% changed nothing for rings of the simulated size. These results are conditional on the simulator: the ring moves the price by about 0.5\% and never reaches the band, whereas the case stocks reached the ceiling for a median of five days and, for one stock, 12 in a row (the band of that stock is assumed) (Table~\ref{tab:casestats}). Whether the band constrains a ring that ratchets the price over days is a question for account-level data.

\paragraph{Do not undo the 2016 tick reform.} On the data we have, the reduction of the tick on HOSE in 2016 lowered the share of zero-return days by 7.1 points relative to control exchanges, a larger change than at four placebo dates, and the estimated spread by about 8\% (the interval includes zero), in the direction the simulator predicts and consistent with lower trading costs found on intraday data \citep{vo2023}. Tick size is a liquidity instrument and is not a manipulation instrument: we found no evidence that it deters rings, and the doubling of the tick in the simulator costs honest participants a third of the spread.

Table~\ref{tab:options} puts the options side by side with the strength of the evidence behind each and the first step a regulator could take.

\begin{table}[t]
\centering\footnotesize
\begin{adjustbox}{max width=\linewidth}
\begin{tabular}{p{2.7cm}p{3.3cm}p{2.5cm}p{3.0cm}p{2.9cm}}
\toprule
option & effect found & evidence & cost or risk & first step\\
\midrule
price off-book blocks from a trailing mean ($L\approx$ one day) & removes 87--98\% of the block value; total ring gain falls to 28--44\% (for the assumed block of 3000 lots) & simulated, two families, with an exact ex-post check and a closed-form law & 0.2--0.5\% gap to the last price on block trades (not measured on real trades) & compare negotiated-trade prices with trailing means on a few stocks\\
higher fine multiple, or more inspection capacity & doubling the multiple from 5 to 10 cuts the capacity needed to deter by a factor of 2.4 & simulated; closed form fitted to the agent-based market at about 100 accounts and extrapolated with $N^{-1/2}$; weaker in markets with more accounts (2 of 4 without trending) & legal change; staff; false flags & compute the workload $BSW$; plan capacity by number of accounts per stock\\
screening in shadow mode on account-level data & measures the share of accounts flagged and the effective number of accounts & real order flow elsewhere shows the null fails; none for Vietnam & data governance; interpretation of flags & planted-ring audit on the regulator's own history\\
no minimum resting time, cancel fee or larger tick as anti-manipulation tools & avoids +7--8\% spread, +21--24\% volatility, +35\% spread & simulated; tick reform on real data & none & none\\
change of the daily band & none for rings of the simulated size & simulated; case stocks reach the ceiling & unknown for rings that ratchet the price & study case stocks with account data\\
\bottomrule
\end{tabular}
\end{adjustbox}
\caption{Options for the regulator, the effect found, the strength of the evidence, cost or risk, and a first step.}\label{tab:options}
\end{table}

\subsection{What a regulator should do to learn more before relying on the test}

\paragraph{Run the test in shadow mode on account-level data.} The one thing that would change our conclusions most is access to account-level order data, which only the exchange, the depository and the commission hold. On such data the first deliverable is not a ranking of accounts but the calibration check of Section~\ref{sec:real-accounts}: the share of accounts with small $p$-values by activity tier, which measures how far the exchangeability condition is violated in that market.

\paragraph{Audit the detector with planted rings in the regulator's own history.} The planted-ring experiment on real order flow (Section~\ref{sec:real-accounts}) is a procedure that does not need labels: inject synthetic coordinated orders into a historical window, vary the size, activity and number of accounts, and measure catch@$B$ against the real background of the stock in question. It yields the capacity $B$ needed per stock for a target detection probability and reveals stocks in which the background of automated accounts makes the test uninformative. Because it uses the venue's own flow, it also estimates $s_0$ of~\eqref{eq:fit} for that venue.

\paragraph{Treat a flag as a screening signal, not as evidence.} The test controls false alarms per account only under an exchangeability condition that real order flow violates (Section~\ref{sec:real-accounts}); with ranking the guarantee is lost (Proposition~\ref{prop:budget}); and in an algorithm-dominated venue the highest-ranked accounts are likely to be honest bots that react to the same signals. The ranking is therefore an input to human review, with a documented route to appeal, not a basis for sanctions. In the simulator an honest fund that splits orders over sub-accounts is flagged at six times the nominal rate, and a regulator that acts on flags must expect to meet such cases.

\subsection{Risks and unintended effects}

\paragraph{Who bears the cost.} The block-price rule falls on those who trade off the order book; in the simulator it has no effect on honest trading because honest agents do not trade blocks, and its cost on real negotiated trades is not measured. Order-level frictions raise volatility and spreads for all participants (Table~\ref{tab:matrix}); the simulated retail trading cost falls under larger ticks and fees because simulated noise traders mostly post passive orders, which we regard as an artefact, so we do not rank who pays most. The fine falls on manipulators. Inspection has a burden of its own: with $B=3$ inspected accounts in each of 600 active accounts per window and 200 windows a year, an account chosen at random would be inspected $3\cdot200/600=1$ time a year, so that every honest account faces about one inspection a year if the ranking were uninformative, and in practice more for loud and automated accounts, which the ranking favours (Section~\ref{sec:real-accounts}). The burden is small in effort and potentially large in reputation, which is a reason to keep the first inspection a data request and not a sanction.

\paragraph{Displacement.} Under inspection the best-responding ring does not stop; it shrinks, to about 3\% of its unconstrained gain in the stock-level family. Deterrence of large rings may push manipulation toward smaller operations that fall below the inspection threshold, which is the right objective if harm is proportional to size and the wrong one if many small rings together do equal harm. The law of Proposition~\ref{prop:power} says where the threshold lies: a ring of $k$ co-active accounts is detectable if $k-1\gtrsim c\sqrt N$, so the threshold rises with the square root of the number of accounts.

\paragraph{Evasion once the rule is known.} A published rule invites a response. Our strategy space includes jitter (spreading members' actions in time) and cross trades, but the ring cannot, for example, trade across several stocks, split its activity over several windows with a different group of accounts, or imitate the order flow of an honest group. Randomising the inspected set among the top-ranked accounts, and keeping the exact statistic confidential, are measures we have not tested.

\paragraph{Privacy and governance.} Account-level surveillance data are sensitive. A shadow-mode pilot should run inside the exchange or the commission, with the output reduced to ranked pseudonymous identifiers, and with the false-flag statistics of the pilot published before any use for enforcement.

\subsection{Generalisation}

The structure that drives the results is not specific to Vietnam: a market dominated by retail investors, a daily price band, off-book or negotiated trades priced from a reference, administrative fines that are a multiple of the gain, and a regulator with limited inspection capacity. For another market, the quantities to re-estimate are the cancel share, the retail share, the volatility of the index and of a typical stock, the number of active accounts per stock, and the reference-price rule for blocks; the family-of-markets design then gives the distribution of the ring's value over markets that match those moments. The introduction of a new trading system in Vietnam in May 2025 may raise the share of automated accounts; Section~\ref{sec:real-accounts} shows that on a venue dominated by automated accounts the coincidence test flags many accounts and loses power, which is a reason to calibrate it on the new flow before relying on it; we did not vary the share of automated accounts.

\subsection{Limitations}

No real account-level data from Vietnam have been used, so we do not know whether real rings are separable from honest coordination by this test. The simulator's manipulators and honest groups were written by the same author as the detectors, and a model trained with labels in the simulator beats the label-free test in every scenario. The calibrated markets do not reproduce the spread, the tails or the absence of serial correlation in daily returns (Figure~\ref{fig:calibration}; a third family that restores it is analysed in Section~\ref{sec:acf}). Rule facts come from broker summaries and press, and we could not establish the pre-reform tick. The ring search has a small budget and the strategy set was found in the first market; the strategies were not re-optimised in each market of the families, so the reported effects of the rules are lower bounds on what an adaptive ring gains (Proposition~\ref{prop:lower}). The block value of 3000 lots is an assumption; a ring that is caught is assumed to be exposed by one inspected member; and fixed manipulators other than the ring are not best responses. The power law is an approximation fitted to 11 points with two constants. The real account-level data are pseudonymous addresses on a crypto derivatives venue, and the planted ring is synthetic. The Vietnamese checks are at the level of the stock; the case periods come from press reports; the panel holds only stocks that are still listed; the natural experiment compares HOSE with other exchanges in the same weeks and cannot rule out other differences.

\subsection{Next steps}

Account-level data from the exchange, or the per-account trade tables in published case files, would allow a direct test of the central claim. A rule based on finer resampling and a false-discovery-rate bound should replace top-$B$ ranking, whose sensitivity falls with the number of accounts. A test conditioned on the common drivers of order flow (price moves, news, activity regimes) is needed for algorithm-dominated markets. The interplay of an adaptive ring with an adaptive regulator is a repeated game that double-oracle methods could solve \citep{mcmahan2003,lanctot2017}. And the optimal combination of fine, capacity and rules is a design problem that the closed forms of Section~\ref{sec:deter} make tractable.

\section{Conclusion}

A regulator without labelled cases can screen for coordinated accounts with a test whose false-alarm rate is exact per account under an exchangeability condition, whose power falls like the inverse square root of the number of accounts, and whose failure profile we have measured on simulated honest groups and on real order flow. Against a ring that adapts, pricing off-book blocks at a trailing mean and a sufficient inspection capacity under a fine proportional to the gain remove most of the ring's expected benefit in simulated markets that match the observable moments of the Ho Chi Minh Stock Exchange (for the block rule, when the block is the main source of the ring's gain), while the daily band does not, and minimum resting times and cancel fees, tested in one market against fixed manipulators, do not and cost honest investors liquidity. The central uncertainty is the one data would settle: whether real rings in a retail-dominated market look like the simulated ones against a background of honest accounts that is as synchronised as the real one.
